% А1.2-1 — Модуль как расстояние: |x − 2| ≤ 3
\documentclass[tikz,border=3pt]{standalone}
% Шапка всех рисунков: те же шрифты и стиль, что в книге.
\def\pfFontPath{../../shrifty/}
\usepackage{fontspec}
\usepackage{unicode-math}
\setmainfont{FiraSans}[Path=\pfFontPath, Extension=.otf,
  UprightFont=*-Regular, ItalicFont=*-Italic, BoldFont=*-SemiBold, BoldItalicFont=*-BoldItalic]
\setmathfont{FiraMath-Regular}[Path=\pfFontPath, Extension=.otf]
\usepackage{stil/profig}

\begin{document}
\begin{tikzpicture}[pf, x=7mm]
  % числовая прямая
  \draw[given, ->] (-3,0) -- (7.3,0) node[below left=1pt] {$x$};
  \foreach \t in {-2,...,7} \draw[pfBlack, line width=\pfThin] (\t,-2pt) -- (\t,2pt);
  % искомый отрезок [-1; 5]
  \draw[find] (-1,0) -- (5,0);
  \fill[pfOrange] (-1,0) circle[radius=2.2pt] (5,0) circle[radius=2.2pt];
  \node[side, below=3pt] at (-1,0) {$-1$};
  \node[side, below=3pt] at (5,0) {$5$};
  % центр и расстояния
  \fill[pfBlue] (2,0) circle[radius=2pt];
  \node[side, below=3pt, text=pfBlue] at (2,0) {$2$};
  \draw[main, <->] (2,0.45) -- node[side, above, text=pfBlue] {$3$} (5,0.45);
  \draw[main, <->] (-1,0.45) -- node[side, above, text=pfBlue] {$3$} (2,0.45);
\end{tikzpicture}
\end{document}
